\section{Algoritmo Planteado y Complejidad}

\subsection{Implementación}
Para resolver el problema planteado, se propone un enfoque Greedy priorizando a los ayudantes que más tiempo requieren. A continuación se presenta el código de la solución implementado en el lenguaje Python:

\begin{lstlisting}[language=Python, caption={Algoritmo Greedy para la planificación de análisis de videos}, frame=single, backgroundcolor=\color{gray!5}]
def planificar_videos(videos):
    videos_ordenados = sorted(videos, key=lambda x: x[2], reverse=True)

    tiempo_scaloni = 0
    tiempo_total = 0

    for s_i, a_i in videos_ordenados:
        tiempo_scaloni += s_i
        tiempo_fin_ayudante = tiempo_scaloni + a_i
        tiempo_total = max(tiempo_fin_ayudante, tiempo_total)

    return videos_ordenados, tiempo_total
\end{lstlisting}

\subsection{Análisis de Complejidad}

Para analizar la complejidad temporal del algoritmo propuesto, desglosaremos su comportamiento en dos fases principales, definiendo a $n$ como la cantidad total de videos (rivales):

\begin{enumerate}
    \item \textbf{Fase de ordenamiento:} La primera operación del algoritmo consiste en ordenar la lista de tuplas \texttt{videos} en función del tiempo de los ayudantes ($a_i$). En Python, la función nativa \texttt{sorted()} utiliza el algoritmo Timsort (un algoritmo de ordenamiento por comparación derivado de Merge Sort e Insertion Sort), cuyo peor caso de tiempo de ejecución es estrictamente $\mathcal{O}(n \log n)$.

    \item \textbf{Fase de iteración (Selección Greedy):} Se inicializan dos variables numéricas simples, lo cual toma tiempo constante $\mathcal{O}(1)$. Luego, el bucle \texttt{for} recorre la lista ordenada exactamente una vez. Dentro del bucle, todas las operaciones ejecutadas (sumas, asignación de variables y la evaluación de la función \texttt{max}) son operaciones elementales que se resuelven en tiempo constante $\mathcal{O}(1)$. Al iterar sobre los $n$ elementos, el tiempo total de esta fase es estrictamente $\mathcal{O}(n)$.
\end{enumerate}

\textbf{Complejidad Temporal Final:}
La complejidad total del algoritmo es la suma de los costos de sus fases:
\[ T(n) = \mathcal{O}(n \log n) + \mathcal{O}(n) \]

Por las propiedades de la notación asintótica (específicamente la regla de la suma), para valores de $n$ suficientemente grandes, los términos de menor orden y los factores constantes son dominados por el término de mayor tasa de crecimiento. Como la función $n \log n$ domina estrictamente a la función lineal $n$, podemos descartar el término $\mathcal{O}(n)$.

Por lo tanto, la complejidad temporal en el peor de los casos para este algoritmo es de \textbf{$\mathcal{O}(n \log n)$}.

\textbf{Complejidad Espacial:}
En términos de espacio, la función \texttt{sorted()} de Python crea y devuelve una nueva lista, lo que requiere espacio adicional directamente proporcional a la cantidad de elementos. Las variables acumuladoras utilizan espacio $\mathcal{O}(1)$. Por consiguiente, la complejidad espacial auxiliar de este algoritmo es $\mathcal{O}(n)$.

\subsection{Resultados del algoritmo}

Se comparan los resultados obtenidos sobre los datos que nos provee la cátedra (resultados óptimos verificados). A continuación, se detalla el orden de los videos procesados y los tiempos de finalización logrados frente a los óptimos teóricos (exluyendo el orden del archivo 10000 elem.txt).

\begin{lstlisting}[frame=tb, numbers=left, basicstyle=\ttfamily\small, breaklines=true, breakindent=0pt]
Archivo: 3 elem.txt
Orden de videos: [(1, 8), (3, 3), (5, 1)]
OPTIMO. Esperado: 10, Obtenido: 10

Archivo: 10 elem.txt
Orden de videos: [(2, 9), (4, 8), (1, 6), (1, 5), (6, 5), (2, 4), (1, 3), (4, 3), (2, 2), (5, 1)]
OPTIMO. Esperado: 29, Obtenido: 29

Archivo: 100 elem.txt
Orden de videos: [(66, 100), (69, 99), (44, 97), (16, 97), (57, 95), (25, 94), (34, 93), (31, 92), (99, 91), (42, 90), (63, 90), (34, 89), (21, 89), (99, 86), (10, 85), (82, 83), (57, 83), (41, 83), (44, 82), (16, 81), (35, 80), (62, 79), (54, 79), (19, 79), (43, 78), (44, 76), (64, 73), (62, 72), (35, 72), (37, 72), (32, 70), (86, 70), (23, 70), (21, 69), (22, 68), (78, 66), (58, 65), (90, 61), (85, 61), (25, 61), (13, 61), (50, 61), (41, 60), (78, 60), (51, 60), (61, 59), (96, 58), (41, 57), (91, 55), (70, 54), (86, 53), (29, 53), (56, 52), (53, 51), (75, 49), (83, 47), (67, 47), (26, 47), (37, 47), (92, 46), (16, 45), (72, 45), (72, 45), (23, 45), (13, 44), (38, 43), (17, 42), (80, 40), (51, 40), (58, 39), (47, 39), (91, 38), (36, 37), (16, 37), (86, 36), (62, 35), (82, 34), (69, 34), (58, 32), (15, 31), (97, 29), (100, 27), (86, 23), (65, 22), (83, 22), (17, 21), (20, 20), (16, 19), (42, 18), (70, 18), (43, 17), (70, 16), (34, 16), (37, 15), (43, 15), (41, 14), (76, 14), (72, 13), (33, 12), (45, 10)]
OPTIMO. Esperado: 5223, Obtenido: 5223

Archivo: 10000 elem.txt
OPTIMO. Esperado: 497886735, Obtenido: 497886735
\end{lstlisting}


